\begin{lemma}[The reversed Kempe path is a Tashkinov initial tree]
Let $(G,e_0,\varphi)$ be a $k$-triple and let $T$ be a tree-sequence with
first two vertices $y_0,y_1$.  Suppose
$\alpha\in\overline\varphi(y_0)$ and
$\beta\in\overline\varphi(y_1)$.  If $P$ is a simple incidence path from
$y_0$ to $y_1$ in the $(\alpha,\beta)$-incidence graph and
$4\mid\operatorname{length}(P)$, then $P$ admits a reversed Kempe-path
certificate in the sense of \noderef{ScheideReversedPathCertificate}.
\end{lemma}
\begin{proof}
By \noderef{TwoColorIncidenceGraph}, every adjacency joins a vertex-node
to an incident edge-node colored $\alpha$ or $\beta$.  Hence, starting at
the vertex-node $y_0$ and ending at the vertex-node $y_1$, write
\[
 P=(\mathrm{inl}(a_0),\mathrm{inr}(f_1),\mathrm{inl}(a_1),
       \ldots,\mathrm{inr}(f_m),\mathrm{inl}(a_m)),
\]
where $a_0=y_0$, $a_m=y_1$, and
$2m=\operatorname{length}(P)$.  Simplicity says that the $a_r$ are
pairwise distinct and that the individual identities $f_r$ are pairwise
distinct.  This does not identify parallel edges.

For $1\leq r\leq m$, the two walk adjacencies at $f_r$ give, by
\noderef{TwoColorIncidenceGraph}, incidence of $f_r$ with both
$a_{r-1}$ and $a_r$, as well as
$\varphi(f_r)\in\{\alpha,\beta\}$.  The consecutive vertex-nodes are
distinct by simplicity.  Expanding incidence with \noderef{Incident} and
using looplessness from \noderef{LooplessMultigraph}, the two distinct
incident vertices are the two endpoints, so
\[
             G.\operatorname{ends}(f_r)=\{a_{r-1},a_r\}.       \tag{1}
\]
The endpoints $y_0,y_1$ are distinct by vertex injectivity in
\noderef{TreeSequence}; therefore $P$ has positive length.  Since its
length is divisible by four, $m\geq2$.  Moreover no $f_r$ equals $e_0$:
each $f_r$ is colored, whereas \noderef{KTriple} and
\noderef{IsOneEdgeUncolored} say that $e_0$ is uncolored.

Define a length-$m$ sequence $U$ by the exact zero-based entries
\[
 \begin{array}{c|c@{\qquad}c|c}
 j&U.vertex(j)&j&U.edge(j)\\ \hline
 0&a_0&0&e_0\\
 1&a_m&1\leq j\leq m-1&f_{m+1-j}\\
 2\leq j\leq m&a_{m+1-j}&&
 \end{array}                                                     \tag{2}
\]
and put $q=f_1$.  The displayed vertex map and edge map are injective by
the preceding pairwise-distinctness statements and $f_r\ne e_0$.
The first edge is $e_0$.  The attachment clause of the original
tree-sequence at index zero says that $e_0$ joins $y_1$ and $y_0$, hence
it attaches the second vertex of $U$ to the first.  If
$1\leq j\leq m-1$, put $r=m+1-j$.  Equation (1) and (2) give
\[
 G.\operatorname{ends}(U.edge(j))
 =G.\operatorname{ends}(f_r)
 =\{a_{r-1},a_r\}
 =\{U.vertex(j),U.vertex(j+1)\}.
\]
Thus every remaining edge attaches its new vertex to the immediately
preceding vertex.  This verifies all fields of \noderef{TreeSequence}.

To verify \noderef{TashkinovTree}, consider an index $j\geq1$.
The edge $U.edge(j)=f_{m+1-j}$ is colored $\alpha$ or $\beta$.  In the
first case its earlier missing-color witness is $U.vertex(0)=y_0$, by the
hypothesis on $\alpha$; in the second it is
$U.vertex(1)=y_1$, by the hypothesis on $\beta$.  Both witnesses precede
the new vertex at index $j+1$.  There is no Tashkinov condition at edge
index zero.

Finally, (2) is a permutation of the vertex-nodes
$a_0,\ldots,a_m$.  Consequently those nodes are exactly the vertex support
of $P$.  Its edge support is exactly $f_1=q$ together with
$f_m,f_{m-1},\ldots,f_2$, the edges of $U$ at indices $1,\ldots,m-1$.
Equation (1) gives the closing endpoint identity for $q$ and every stated
trunk endpoint identity.  Together with $m=\operatorname{length}(P)/2$
and $m\geq2$, these are precisely all fields of the reversed Kempe-path
certificate.
\end{proof}
